% Pursuing Stacks (cotes 134-2 to 134-8): archive cote 134-2, site page 49 (PDF page 50).
% appendix (letters to Breen). Draft compared: ps1.tex:1245-1252.
% AI-transcribed by Claude Opus 5.5 (claude-opus-5-5, adaptive thinking, effort high), Message Batches API, batch ps1 (msgbatch_01Rv5hCaNTtUt4JHx1yqoHPL), 2026-11-12; not yet checked by eye.
\scanpage{134-2/49}
\mspage{39} \hfill 2

son "(pro)-type d'homotopie n-tronquée", et le définissent.  Dans le  cas
où  X est localement homotopiquement trivial en dim \hand{$\leq$} n, donc définit un
type d'homotopie  n- tronqué ordinaire, on peut interpréter ce dernier comme
\margin{$\ulcorner$ \unsure{(C)} $\llcorner$} un n-groupoide  $\hand{\mathcal{C}}_{n}$, (défini à n-équivalence près).\hand{$\|$}En termes de \deleted{\unsure{ceux-ci}}\inserted{$\mathcal{C}_n$}\hand{,}
les (n-1) - champs \inserted{localement constants} sur X doivent s'identifier aux n-foncteurs de la n- caté-%
gorie  $\hand{\mathcal{C}}_{n}$  dans  la n-catégorie   ( (n-1)- Cat)  de toutes les (n-1) - caté-%
gories.\hand{$\|$}Dans le cas n=1  ceci n'est autre que la théorie  de Poincaré  de la
classification des revêtements de X en termes du "groupoide fondamental" $\hand{\mathcal{C}}_{1}$
de X.  Par extension, $\hand{\mathcal{C}}_{n}$  mérite le nom de  \uline{n-groupoide fondamental  de}
X, que je propose de noter $\hand{\Pi}_{n}(X)$.  Sa connaissance \unsure{induite|induit} donc celle
des  $\hand{\pi}_{i}(X)$  ($0 \leq i \leq n$)  et des invariants de Postnikoff  de tous les
ordres jusqu'à  $H^{n+1}$  ($\hand{\Pi}_{n-1}(X)$ , $\hand{\pi}_{n}$).

Dans le cas d'un topos X quelconque, pas \unsure{récessairement|nécessairement} \inserted{loc.} homotopique-%
ment trivial en dim $\leq$ n,  on espère pouvoir interpréter les (n-1) - champs
localement constants sur X   enttermes d'un $\hand{\Pi}_{n}(X)$  qui sera un  pro-n-groupoide.
Ça  a été fait en tous cas, plus ou moins, pour n = 1  (du moins pour X conne-%
xe);  le cas où  X est le topos étale d'un schéma  est traité  in extenso
dans SGA 3, à propos de la classification des tores sur une base quelconque.

Dans le cas n = 1,  on sait qu'on récupère (à équivalence près)   le
1-groupoide  $\hand{\mathcal{C}}_{1}$  à partir de la 1- catégorie  \uline{Hom} ($\hand{\mathcal{C}}_{1}$, (Ens))  de ces
foncteurs  dans (Ens) = (0-cat) (i.e. des "systèmes locaux" sur $\hand{\mathcal{C}}_{1}$ )\deleted{\illegible{1}}\hand{(}qui est
un  topos , dit "multigaloisien"\hand{)}\deleted{\unsure{,}}  comme la catégorie des "foncteurs fibres "
sur ledit topos, i.e. la catégorie opposée à la catégorie des points de ce
topos (\deleted{\unsure{qui}}\inserted{lequel} n'est autre que le  \uline{topos classifiant}  de $\hand{\mathcal{C}}_{1}$). Pour préciser
pour n quelconque la façon dont le n-type d'homotopie d'un topos X (supposé
localement homotopiquement trivial en dim $\leq$ n, pour simplifier), i.e.  son
n- groupoide fondamental $\hand{\mathcal{C}}_{n}$,  s'exprime en termes de la  n-catégorie des

\mspage{49}
